% useful command to build & clean-up : 
% latexmk -pdf -outdir="French/Portfolio/PDF" -c "French/Portfolio/dossier_competences.tex"
% latexmk -pdf -outdir="French/Portfolio/PDF" "French/Portfolio/dossier_competences.tex"
% latexmk -c "French/Portfolio/dossier_competences.tex"
% find French/Portfolio/PDF -type f \( -name "*.aux" -o -name "*.log"  \) -delete


\documentclass{article}
\usepackage[scaled]{helvet}
\usepackage{amssymb}
\usepackage{textcomp} % for \textquotesingle
\renewcommand{\familydefault}{\sfdefault} % applies font
\usepackage{tcolorbox}
\usepackage{xcolor}

\usepackage[left=0.1in, right=0.1in, top=0.55in, bottom=0.55in]{geometry}
\usepackage{graphicx}
\usepackage{tikz}
\usepackage{hyperref}
\usepackage{enumitem}
\usepackage{fancyhdr}
\usepackage{titlesec}
\usepackage{microtype}
\usepackage{colortbl}
\usepackage{array} % Required for m{...} column type
\usepackage{lastpage} % Required for \pageref{LastPage} to get the total number of pages




% Page style setup

\definecolor{titleBlue}{RGB}{0, 102, 204}
\definecolor{skyblue}{RGB}{135,206,235}
\definecolor{deepBlue}{RGB}{80,100,255} 
\definecolor{mainBlue}{RGB}{0, 51, 102} 
\definecolor{secondaryBlue}{RGB}{0, 153, 204} %
\definecolor{steelBlue}{RGB}{70, 130, 180}
\definecolor{darkCyan}{RGB}{0, 139, 139}
\definecolor{cream}{RGB}{255,253,208}
\definecolor{darkGray}{RGB}{54, 69, 79}
\definecolor{deepPurple}{RGB}{121, 53, 149}
\definecolor{enseaRed}{RGB}{191, 0, 80}
\definecolor{antiEnsea}{RGB}{32, 127, 87}
\definecolor{professionalOrange}{RGB}{200, 100, 0}

% Custom command to display current page number and total number of pages
\newcommand{\PageNumbering}{%
    \ifnum\value{page}>1%
        \thepage\ / \pageref{LastPage}%
    \else%
        \thepage%
    \fi%
}

\pagestyle{fancy}
\fancyhf{}
\fancyhead[L]{\textcolor{darkGray}{\textbf{Dossier de compétences}}} % Left header
\fancyhead[R]{\textcolor{darkGray}{\textbf{Alfred LALANNE}}} % Right header
\fancyfoot[R]{\PageNumbering} % Right footer with custom page numbering
\renewcommand{\headrulewidth}{1pt} % Header rule width
\renewcommand{\headrule}{\color{mainBlue}\hrule width\headwidth height\headrulewidth \vskip-\headrulewidth} % Redefine color of header rule
\title{Ingénieur Électronique et Logiciel}


\date{} % Remove the date

% Redefine \maketitle to remove the date and include photo and colored title
\makeatletter
\renewcommand{\maketitle}{%
    \begin{center} % Center the whole content    
        {\Large\bfseries{\textcolor{darkGray}{\@title}} \par}
        \begin{minipage}[b]{0.6\linewidth} % left (empty) minipage
        \end{minipage}%
        \hfill % Add some horizontal space between the minipages
        \begin{minipage}[b]{0.3\linewidth}   % right minipage
            \raggedleft
            \vspace*{-1.78\baselineskip} % Adjust the vertical position of the image
            \includegraphics[width=0.5\linewidth]{moi.jpeg} 
            \hspace{0.1\linewidth}
        \end{minipage}%
    \end{center} % End of center environment
}
\makeatother


\newcolumntype{C}[1]{>{\centering\arraybackslash}m{#1}}

\begin{document}


\maketitle

%\section{Domaines de compétence}



%%%%%%%%%%%%%%%%%%%%%%%%%%%%%%%%%%%%%%%%%%%%%%%%%%%%%% DOMAINES DE COMPETENCES %%%%%%%%%%%%%%%%%%%%%%%%%%%%%%%%%%%%%%%%%%%%%%%%%%%%%%%%%%%%%%%%%%%%%%%

\begin{center}
    \begin{tikzpicture}
        \draw[mainBlue, line width=1pt] (-0.4\linewidth,-2pt) -- (0.4\linewidth,-2pt);
    \end{tikzpicture}
\end{center}

\begin{tabular}{@{\hspace{0.05\textwidth}}m{0.1\textwidth}@{\hspace{0.05\textwidth}}!{\color{secondaryBlue}\vline width 2pt}@{}m{0.7\textwidth}@{}}
    \textcolor{secondaryBlue}{\textbf{Domaines de compétences}} &
    \begin{itemize}[label={}, topsep=0pt, partopsep=0pt, itemsep=0pt, parsep=0pt]
        \setlength{\itemsep}{0pt} 

        \item \textcolor{darkGray}{Hardware - Electronique}
        \begin{itemize}[label={\textcolor{gray!80}{\checkmark}}, topsep=0pt, partopsep=0pt, itemsep=0pt, parsep=0pt] 
            \item \textcolor{gray!80}{Conception de schémas électriques, routage, simulation }
            \item \textcolor{gray!80}{Filtrage analogique actif et passif, Matlab}
        \end{itemize}

        \item \textcolor{darkGray}{Electronique numérique} 
        \begin{itemize}[label={\textcolor{gray!80}{\checkmark}}, topsep=0pt, partopsep=0pt, itemsep=0pt, parsep=0pt] 
            \item \textcolor{gray!80}{Microcontrôleurs : programmation en C et Assembleur sur Keil-uVision }
            \item \textcolor{gray!80}{Logique programmable : technologie FPGA, description VHDL}
            \item \textcolor{gray!80}{Systèmes sur puce : SoC, SoC-FPGA, SoC-IP}
        \end{itemize}

        \item \textcolor{darkGray}{Informatique}
        \begin{itemize}[label={\textcolor{gray!80}{\checkmark}}, topsep=0pt, partopsep=0pt, itemsep=0pt, parsep=0pt] 
            \item \textcolor{gray!80}{Programmation bas-niveau : microcontrôleurs, FreeRTOS}
            \item \textcolor{gray!80}{Informatique industrielle : traitement d'image, création d'interfaces graphiques}
            \item \textcolor{gray!80}{Développement jeux-vidéos}
        \end{itemize}

        \item \textcolor{darkGray}{Ingénierie système}
        \begin{itemize}[label={\textcolor{gray!80}{\checkmark}}, topsep=0pt, partopsep=0pt, itemsep=0pt, parsep=0pt, after=\vspace*{-\baselineskip}] % reduce spacing & adjust left margin
            \item \textcolor{gray!80}{Rédaction de cahier des charges}
            \item \textcolor{gray!80}{Etudes de faisabilité}
        \end{itemize}
    \end{itemize}
\end{tabular}


%%%%%%%%%%%%%%%%%%%%%%%%%%%%%%%%%%%%%%%%%%%%%%%%%%%%%% LANGAGES OUTILS NORMES %%%%%%%%%%%%%%%%%%%%%%%%%%%%%%%%%%%%%%%%%%%%%%%%%%%%%%%%%%%%%%%%%%%%%%%

\begin{center}
    \begin{tikzpicture}
        \draw[mainBlue, line width=1pt] (-0.4\linewidth,-2pt) -- (0.4\linewidth,-2pt);
    \end{tikzpicture}
\end{center}


\begin{tabular}{@{\hspace{0.05\textwidth}}m{0.1\textwidth}@{\hspace{0.05\textwidth}}!{\color{secondaryBlue}\vline width 2pt}@{}m{0.7\textwidth}@{}}
    \textcolor{secondaryBlue}{\textbf{Langages\newline Outils\newline Normes}} & 
    \begin{itemize}[label={}, topsep=0pt, partopsep=0pt, itemsep=0pt, parsep=0pt] 
        \setlength{\itemsep}{0pt} % Reduce spacing between lines

        \item \textcolor{darkGray}{Langages}
        \begin{itemize}[label={\textcolor{gray!80}{\checkmark}}, topsep=0pt, partopsep=0pt, itemsep=0pt, parsep=0pt] 
            \item \textcolor{gray!80}{C, C++, Python, Java, Assembleur}
        \end{itemize}

        \item \textcolor{darkGray}{Outils}
        \begin{itemize}[label={\textcolor{gray!80}{\checkmark}}, topsep=0pt, partopsep=0pt, itemsep=0pt, parsep=0pt]
            \item \textcolor{gray!80}{Collaboratifs : Git, Jira, Confluence, Teams}
            \item \textcolor{gray!80}{IDEs : VSCode, STM32CubeIDE, Eclipse}
            \item \textcolor{gray!80}{Bibliothèques : OpenCV, PyQt5, NumPy, SFML, Matplotlib, RealSense}
            \item \textcolor{gray!80}{Modélisation : OrCAD PSpice, Visio, Blender, draw.io}
        \end{itemize}

        \item \textcolor{darkGray}{Normes}
        \begin{itemize}[label={\textcolor{gray!80}{\checkmark}}, topsep=0pt, partopsep=0pt, itemsep=0pt, parsep=0pt, after=\vspace*{-\baselineskip}] 
            \item \textcolor{gray!80}{ISO7816 (smart-cards), ISO12233 (traitement d'image)}
        \end{itemize}
    \end{itemize}
\end{tabular}


%%%%%%%%%%%%%%%%%%%%%%%%%%%%%%%%%%%%%%%%%%%%%%%%%%%%%% SECTEURS D'ACTIVITES %%%%%%%%%%%%%%%%%%%%%%%%%%%%%%%%%%%%%%%%%%%%%%%%%%%%%%%%%%%%%%%%%%%%%%%

\begin{center}
    \begin{tikzpicture}
        \draw[mainBlue, line width=1pt] (-0.4\linewidth,-2pt) -- (0.4\linewidth,-2pt);
    \end{tikzpicture}
\end{center}


\begin{tabular}{@{\hspace{0.05\textwidth}}m{0.1\textwidth}@{\hspace{0.05\textwidth}}!{\color{secondaryBlue}\vline width 2pt}@{}C{0.7\textwidth}@{}}
    \textcolor{secondaryBlue}{\textbf{Secteurs d'activités}} & 
    \begin{itemize}[label={}, topsep=0pt, partopsep=0pt, itemsep=0pt, parsep=0pt, after=\vspace*{-\baselineskip}]
        \item \textcolor{gray!80}{Systèmes d'identification et de sécurité}
        \item \textcolor{gray!80}{Systèmes d'acquisition automatisés}
        \item \textcolor{gray!80}{Smart-Cards}
    \end{itemize}
\end{tabular}


%%%%%%%%%%%%%%%%%%%%%%%%%%%%%%%%%%%%%%%%%%%%%%%%%%%%%% FORMATION %%%%%%%%%%%%%%%%%%%%%%%%%%%%%%%%%%%%%%%%%%%%%%%%%%%%%%%%%%%%%%%%%%%%%%%

\begin{center}
    \begin{tikzpicture}
        \draw[mainBlue, line width=1pt] (-0.4\linewidth,-2pt) -- (0.4\linewidth,-2pt);
    \end{tikzpicture}
\end{center}

%\begin{tabular}{@{\hspace{0.05\textwidth}}m{0.1\textwidth}@{\hspace{0.05\textwidth}}!{\color{secondaryBlue}\vline width 2pt}@{}m{0.7\textwidth}@{}}

\begin{tabular}{@{\hspace{0.05\textwidth}} m{0.1\textwidth} @ {\hspace{0.05\textwidth}} ! {\color{secondaryBlue}\vline width 2pt \hspace{0.04\textwidth}} @ {}m{0.4\textwidth} @{\hspace{0.15\textwidth}} m{0.5\linewidth}@{}}
    \textcolor{secondaryBlue}{\textbf{Formation}} & 
    \textcolor{darkGray}{Ingénieur électronicien ENSEA}
    \begin{itemize}[label={\textcolor{gray!80}{\checkmark}}, topsep=0pt, partopsep=0pt, itemsep=0pt, parsep=0pt, after=\vspace*{-\baselineskip}]
        \item \textcolor{gray!80}{Année d'obtention du diplôme : 2023} 
        \item \textcolor{gray!80}{Spécialités : microélectronique et numérique} 
    \end{itemize} &
    \includegraphics[width=0.1\linewidth]{ensea.jpg}
\end{tabular}


%%%%%%%%%%%%%%%%%%%%%%%%%%%%%%%%%%%%%%%%%%%%%%%%%%%%%%%%%%%%%%%%%%%%%%%%%%%%%%%%%%%%%%%%%%%%%%%%%%%%%%%%%%%%%%%%%%%%%%%%%%%%%%%%%%%%%%%%%%%%%%%%%%%%





\end{document}